\section{Actual canonical modules and counits on the polynomial Kummer chart}
\label{sec:kn-nadel-kummer-canonical-trace}

Let $R$ be a nontrivial commutative ring and let
$A=R[x_0,\ldots,x_{n-1}]$, $B=R[t_0,\ldots,t_{n-1}]$.
The actual map is $f:A\to B$, $x_j\mapsto t_j^{d_j+1}$.
These two rings have the same underlying polynomial type in Lean,
so all coefficient restriction and coextension objects below use
the literal homomorphism $f$, rather than the identity scalar action.
Let $C$ be the internally constructed coordinate root tower and
$e:C\simeq_A B$ the already proved original coordinate algebra
isomorphism.  Put
$\Lambda(b)=\Lambda_C(e^{-1}(b))$, with $\Lambda_C$ the actual
composite of the original monic top-coefficient functionals.

\begin{lemma}[Transport the actual whole dual frame to the power-map algebra]
\label{lem:kn-nadel-kummer-polynomial-dual}
\lean{EnvelopeProperty.Vanishing.NadelKummerCanonicalTrace.coextensionFrame,EnvelopeProperty.Vanishing.NadelKummerCanonicalTrace.coextensionFrame_apply,EnvelopeProperty.Vanishing.NadelKummerCanonicalTrace.PolynomialRing,EnvelopeProperty.Vanishing.NadelKummerCanonicalTrace.powerMap,EnvelopeProperty.Vanishing.NadelKummerCanonicalTrace.towerRestrictionEquiv,EnvelopeProperty.Vanishing.NadelKummerCanonicalTrace.powerDualizingModule,EnvelopeProperty.Vanishing.NadelKummerCanonicalTrace.coefficientMap,EnvelopeProperty.Vanishing.NadelKummerCanonicalTrace.coefficientMap_apply,EnvelopeProperty.Vanishing.NadelKummerCanonicalTrace.powerDualizingMap,EnvelopeProperty.Vanishing.NadelKummerCanonicalTrace.powerDualizingMap_apply,EnvelopeProperty.Vanishing.NadelKummerCanonicalTrace.functionalTransport,EnvelopeProperty.Vanishing.NadelKummerCanonicalTrace.functionalTransport_apply,EnvelopeProperty.Vanishing.NadelKummerCanonicalTrace.functionalTransport_powerDualizingMap,EnvelopeProperty.Vanishing.NadelKummerCanonicalTrace.powerDualizingMap_bijective,EnvelopeProperty.Vanishing.NadelKummerCanonicalTrace.powerDualizingEquiv,EnvelopeProperty.Vanishing.NadelKummerCanonicalTrace.powerDualizingIso}
\leanok
Construct a $B$-linear equivalence
\[
 B\simeq\operatorname{Coext}_f(A),\qquad
 g\longmapsto[b\mapsto\Lambda(bg)].
\]
The map $e$ preserves the original coefficient homomorphisms on the
whole original algebra.  No comparison, perfect pairing, or desired
duality isomorphism is an input to this construction.
\end{lemma}
\begin{proof}
Use the actual coefficient triangle of $e$ to produce an $A$-linear
equivalence between the two native restricted regular modules.
Precomposition with this equivalence identifies their entire
native coextension carriers.  Compose the original tower's actual
coefficient counit with this restriction equivalence to obtain
$\Lambda$ as an actual $A$-linear map on the restricted $B$.
The existing native restriction/coextension adjunction sends this
functional to the displayed $B$-linear map.  Under precomposition
with $e$, this is exactly the already proved tower frame evaluated
at $e^{-1}(g)$, by the actual multiplication and inverse identities.
Since both $e$ and the whole tower frame are bijective, the native
map is bijective as well.  Thus it is the asserted $B$-linear
equivalence on the actual power-map algebra.
\uses{thm:kn-nadel-finite-root-tower-duality,lem:kn-nadel-kummer-tower-polynomial-iso}
\end{proof}

\begin{theorem}[The actual top-differential finite-duality map]
\label{thm:kn-nadel-kummer-canonical-module}
\lean{EnvelopeProperty.Vanishing.NadelKummerCanonicalTrace.canonicalModule,EnvelopeProperty.Vanishing.NadelKummerCanonicalTrace.canonicalGenerator,EnvelopeProperty.Vanishing.NadelKummerCanonicalTrace.canonicalTrivialization,EnvelopeProperty.Vanishing.NadelKummerCanonicalTrace.canonicalGenerator_eq_wedge,EnvelopeProperty.Vanishing.NadelKummerCanonicalTrace.canonicalTrivialization_generator,EnvelopeProperty.Vanishing.NadelKummerCanonicalTrace.canonicalTrivialization_symm,EnvelopeProperty.Vanishing.NadelKummerCanonicalTrace.canonicalIso,EnvelopeProperty.Vanishing.NadelKummerCanonicalTrace.canonicalIso_apply,EnvelopeProperty.Vanishing.NadelKummerCanonicalTrace.canonicalIso_generator}
\leanok
Use the genuine K\"ahler modules and top exterior powers
\[
 \omega_A=\bigwedge_A^n\Omega^1_{A/R},\qquad
 \omega_B=\bigwedge_B^n\Omega^1_{B/R}.
\]
Construct on this actual polynomial chart a $B$-linear equivalence
\[
 \omega_B\simeq\operatorname{Coext}_f(\omega_A),\qquad
 g\,dt_0\wedge\cdots\wedge dt_{n-1}
 \longmapsto
 [b\mapsto\Lambda(bg)\,dx_0\wedge\cdots\wedge dx_{n-1}].
 \tag{K1}
\]
Both exterior modules and both displayed generators are the native
original differential objects.  The assertion also includes $n=0$.
\end{theorem}
\begin{proof}
The native polynomial K\"ahler basis is exactly the family of the
original $dX_j$.  Its actual top exterior basis is indexed by the
unique full subset of $\mathrm{Fin}(n)$, giving the displayed
ordered wedge and a coefficient-linear trivialization of each
original $\omega$.  Compose the trivialization of $\omega_B$,
the preceding actual power-map dual frame, and coextension of the
inverse trivialization of $\omega_A$.  Every map has the actual
$B$-action and coefficient homomorphism, hence their composite is
$B$-linear.  Evaluating at every original $b\in B$ gives (K1).
The exterior-basis construction in degree zero is the native
degree-zero exterior module, so the same proof covers that endpoint.
This constructs the map between the genuine differential modules;
it does not rename an abstract dual module as a canonical module.
\uses{lem:kn-nadel-kummer-polynomial-dual,sc0-canonical-port:geom:polynomial-canonical-frame,sc0-canonical-port:geom:canonical-frame-from-differential-basis}
\end{proof}

\begin{corollary}[The original native canonical counit and coefficient formula]
\label{cor:kn-nadel-kummer-canonical-counit}
\lean{EnvelopeProperty.Vanishing.NadelKummerCanonicalTrace.canonicalCounit,EnvelopeProperty.Vanishing.NadelKummerCanonicalTrace.canonicalCounit_apply,EnvelopeProperty.Vanishing.NadelKummerCanonicalTrace.canonicalCounit_generator,EnvelopeProperty.Vanishing.NadelKummerCanonicalTrace.canonicalCounit_composite_topCoefficient}
\leanok
Compose (K1) with the actual native restriction/coextension counit
on $\omega_A$.  The resulting $A$-linear map from the actual
restricted $\omega_B$ to $\omega_A$ sends
\[
 g\,dt_0\wedge\cdots\wedge dt_{n-1}
 \longmapsto\Lambda(g)\,dx_0\wedge\cdots\wedge dx_{n-1}.
 \tag{K2}
\]
\end{corollary}
\begin{proof}
The original native counit evaluates a coextended functional at
$1_B$.  Substitute $b=1$ into (K1) and use multiplication by one.
This proves the entire map formula on every original top form,
because the native top wedge is a basis.
\uses{thm:kn-nadel-kummer-canonical-module}
\end{proof}

The coordinate-normalized construction above is meaningful even
for noninvertible degrees; no field-trace division or etale-locus
argument is used in its construction.  To identify it intrinsically
under arbitrary coordinate changes, further proofs must compare
it with the algebra trace and the actual Jacobian on the etale
locus when the relevant degrees are units, then extend by genuine
regular localization detection.  Etale coefficient base change and
global gluing remain additional producer obligations.  Neither an
ordinary canonical line on a singular normalized cover nor a global
canonical trace or vanishing theorem is asserted here.
